\documentclass[11pt,a4paper]{article}

\usepackage[utf8]{inputenc}
\usepackage[T1]{fontenc}
\usepackage{lmodern}
\usepackage[margin=0.76in]{geometry}
\usepackage[table]{xcolor}
\usepackage{array}
\usepackage{tabularx}
\usepackage{booktabs}
\usepackage{enumitem}
\usepackage{titlesec}
\usepackage{fancyhdr}
\usepackage{lastpage}
\usepackage{hyperref}

\definecolor{TextInk}{HTML}{233142}
\definecolor{HeadingBlue}{HTML}{1D3557}
\definecolor{AccentTeal}{HTML}{2A7F92}
\definecolor{AccentGold}{HTML}{C58B4E}
\definecolor{TableStripe}{HTML}{F3F8FA}
\definecolor{SoftBlue}{HTML}{EEF6F8}
\definecolor{SoftGreen}{HTML}{EEF7F0}
\definecolor{SoftRed}{HTML}{FBF0EF}

\hypersetup{
  colorlinks=true,
  linkcolor=HeadingBlue,
  urlcolor=AccentTeal,
  pdftitle={IB Geography SL May 5 Visual Response Marking Guide},
  pdfauthor={IB Geography SL Preparation}
}

\color{TextInk}
\setlength{\parskip}{0.45em}
\setlength{\parindent}{0pt}
\setlength{\tabcolsep}{5pt}
\renewcommand{\arraystretch}{1.18}
\setlist[itemize]{leftmargin=1.32em,itemsep=3pt,topsep=4pt}

\newcolumntype{Y}{>{\raggedright\arraybackslash}X}
\newcolumntype{L}[1]{>{\raggedright\arraybackslash}p{#1}}

\titleformat{\section}
  {\normalfont\Large\sffamily\bfseries\color{HeadingBlue}}
  {}{0pt}{}
  [\vspace{0.08em}{\color{AccentGold}\titlerule[0.8pt]}]
\titleformat{\subsection}
  {\normalfont\large\sffamily\bfseries\color{AccentTeal}}
  {}{0pt}{}

\pagestyle{fancy}
\setlength{\headheight}{14pt}
\fancyhf{}
\fancyhead[L]{\sffamily\color{AccentTeal}May 5 Marking Guide}
\fancyhead[R]{\sffamily\color{HeadingBlue}IB Geography SL}
\fancyfoot[C]{\sffamily\color{HeadingBlue}\thepage\ of \pageref{LastPage}}
\renewcommand{\headrulewidth}{0.4pt}

\newcommand{\term}[1]{\textcolor{HeadingBlue}{\textbf{#1}}}
\newcommand{\exam}[1]{\textcolor{AccentTeal}{\textbf{#1}}}
\newcommand{\score}[1]{\textcolor{AccentTeal}{\textbf{#1}}}
\newcommand{\boxnote}[3]{%
\noindent\colorbox{#1}{%
  \begin{minipage}{0.965\textwidth}
  \sffamily
  \textbf{\textcolor{HeadingBlue}{#2}} #3
  \end{minipage}
}}

\begin{document}

\begin{center}
  {\sffamily\bfseries\color{HeadingBlue}\LARGE May 5 Visual-Stimulus Marking Guide}\\[0.25em]
  {\sffamily\color{AccentTeal}\large Use After The Timed Response Set Only}
\end{center}

\boxnote{SoftRed}{Strict rule.}{Mark only what was written during the timed
visual-stimulus response. Do not improve the answer before scoring.}

\section{Question 1: Describe Two Patterns Or Contrasts [4]}

Award up to \score{4 marks}: up to \score{2 marks} for each clear pattern or
contrast using Source A evidence.

Valid answers include:

\begin{itemize}
  \item Delta Coast has the highest climate-risk index, at \term{88}, while
  Heatwave City has the lowest climate-risk index, at \term{69}.
  \item Dryland Belt has the highest food-system vulnerability index,
  \term{91}, even though its climate-risk index, \term{76}, is lower than
  Delta Coast's \term{88}.
  \item Heatwave City has the highest health-system stress index,
  \term{84}, but the lowest food-system vulnerability index, \term{44}.
  \item Island Port has high climate risk, \term{83}, and moderate food-system
  vulnerability, \term{68}, and health-system stress, \term{61}.
  \item Delta Coast and Island Port both have high climate-risk scores, but
  Delta Coast has higher food-system vulnerability and health-system stress.
  \item Heatwave City has the highest policy-readiness index, \term{64}, while
  Dryland Belt has the lowest policy-readiness index, \term{35}.
\end{itemize}

\section{Question 2: Different Policy Responses [4]}

Award up to \score{4 marks}: up to \score{2 marks} for each explained reason.
Good answers should link the source to vulnerability, exposure, sensitivity, or
capacity to respond.

\begin{itemize}
  \item Regions may face different dominant hazards: Delta Coast is linked to
  flooding and land-use planning, while Dryland Belt is linked to drought,
  water conservation, and crop adaptation.
  \item Food-system vulnerability differs: Dryland Belt has the highest value,
  \term{91}, so food production, access, prices, or rural livelihoods may need
  stronger policy attention.
  \item Health-system stress differs: Heatwave City has \term{84}, so heat
  plans, healthcare outreach, shade, cooling, and protection for vulnerable
  groups may be more urgent than food-system adaptation.
  \item Policy readiness differs: Dryland Belt has low readiness, \term{35},
  while Heatwave City has higher readiness, \term{64}; implementation capacity
  changes what policies are realistic.
  \item Exposure and sensitivity differ between coasts, drylands, cities, and
  ports, so the same climate-risk index can produce different impacts.
\end{itemize}

\section{Question 3: Policy Response And Limitation [2]}

Award \score{1 mark} for a valid policy response from Source A and
\score{1 mark} for a valid limitation.

Valid policy responses include:

\begin{itemize}
  \item flood shelters, drainage, and land-use planning for Delta Coast;
  \item drought planning, water conservation, or crop adaptation for Dryland
  Belt;
  \item heat-action plans, shade, or healthcare outreach for Heatwave City;
  \item storm-surge planning, port resilience, or import security for Island
  Port.
\end{itemize}

Valid limitations include cost, maintenance, governance capacity, weak policy
readiness, unequal access, difficulty reaching informal or low-income groups,
need for training, energy demand, relocation conflict, or the fact that a
single policy may not address the underlying exposure.

\section{Final Decision}

\rowcolors{2}{TableStripe}{white}
\begin{tabularx}{\textwidth}{L{0.32\textwidth}Y}
\toprule
\textbf{Prompt} & \textbf{Fill in after marking} \\
\midrule
Score out of 10 & \rule{0.65\linewidth}{0.4pt} \\
Main source-use error & \rule{0.65\linewidth}{0.4pt} \\
Main policy-evaluation error & \rule{0.65\linewidth}{0.4pt} \\
One sentence to carry into Paper 2 & \rule{0.65\linewidth}{0.4pt} \\
\bottomrule
\end{tabularx}
\rowcolors{2}{}{}

\boxnote{SoftGreen}{Target.}{A strong May 5 response scores \score{8/10} or
higher by using exact source values, explaining why policy needs differ, and
giving one clear limitation for a response.}

\end{document}
